\chapter*{Preface}
\addcontentsline{toc}{chapter}{Preface}
\markboth{Preface}{Preface}

This book develops the algebra behind the Fourier transform and shows how that
algebra leads to practical algorithms. It is intended for readers who want to
move between mathematics and computation: an engineer seeking to understand why
a transform works, a mathematician looking for applications of finite groups,
or a programmer implementing a spectral method.

\medskip
\noindent\textbf{About this edition.}
\emph{\BookTitle} is an adapted English translation of Gabriel Peyr\'e's
French book \emph{\FrenchBookTitle}, published by Ellipses in 2004
(ISBN 9782729818678). The
\href{\FrenchBookURL}{publisher's page} describes the original French edition.
The present edition revises the exposition, corrects mathematical and
transcription errors, redraws the illustrations, and expresses the computational
examples in Python. The previous English sources and their MATLAB and Maple
programs are preserved in the accompanying \texttt{legacy/} directory.

\medskip
The first two chapters explain Fourier analysis on finite abelian groups and
its applications to duality, quadratic reciprocity, and the Walsh transform.
Chapters~3 and~4 develop the discrete Fourier transform, fast algorithms, and
applications to filtering, interpolation, and differential equations.
Chapter~5 studies related transforms, including the Hartley, $z$-, and
fractional Fourier transforms. Chapter~6 introduces Fourier transforms over
finite fields and their use in error-correcting codes. The final two chapters
extend the theory to noncommutative groups through linear representations and
their characters.

The exposition assumes undergraduate linear algebra and basic familiarity with
groups and complex numbers. Additional prerequisites are introduced where
needed. The chapters on finite fields and representations require more
algebraic maturity. Readers primarily interested in numerical algorithms may
start with Chapter~3 and return to the algebraic foundations as needed.

Exercises connect the theory to explicit calculations and applications. They
are numbered within each chapter, with detailed solutions at the end of that
chapter. The two computational appendices explain
the numerical and exact-arithmetic conventions used in the Python programs.
Numerical examples use NumPy and SciPy, symbolic calculations use SymPy, and
finite-field calculations use the \texttt{galois} package. Displayed programs
are drawn from the accompanying Python source files so that the formulas,
implementations, and executable checks can be studied together.

Convolution is a recurring theme. On a finite group it comes from multiplication
in the group algebra; in signal processing it describes filtering; in
probability it gives the law of a sum of independent variables. The Fourier
transform turns each of these operations into a simpler calculation. Following
that principle across different settings is the organizing thread of the book.
